\section{Advanced Persistent Threat (APT)}
\label{sec:threat-apt}

The attack chain targets the zone aggregator gateways. Phase 0 consists of endpoint enumeration: requests are sent to common
administrative paths in order to discover the honeypot endpoints. Phase 1
consists of repeated honeypot probing to extract the \texttt{INTERNAL\_TOKEN},
which grants the ability to impersonate any internal service. As modelled in the
Attack-Defense Tree (Figure~\ref{fig:attacktree}), both phases must succeed. A
defence that blocks either phase disrupts the full chain. The full attack chain
is demonstrated by \texttt{scripts/demo\_apt.sh}.

\subsection{Detection}

Detection is implemented across four Wazuh rules covering both phases. Rules
100202 and 100203 detect the reconnaissance phase by monitoring for
\texttt{not\_found} events (HTTP 404 responses) on zone aggregator logs. Rule
100203 fires when 5 or more distinct unknown paths are probed from the same IP
within 30 seconds, confirming enumeration behaviour.

\begin{lstlisting}[style=xmlstyle, caption={Wazuh rules for APT recon detection.}]
<rule id="100202" level="3">
  <if_sid>86600</if_sid>
  <field name="event_type">not_found</field>
  <description>Zone aggregator: unknown endpoint
    probed</description>
  <group>apt,recon,</group>
</rule>

<rule id="100203" level="10" frequency="5" timeframe="30">
  <if_matched_sid>100202</if_matched_sid>
  <same_field>client_ip</same_field>
  <description>APT recon: endpoint enumeration -
    5+ unknown paths in 30s from same IP</description>
  <group>apt,recon,attack,</group>
</rule>
\end{lstlisting}

Rules 100200 and 100201 detect the honeypot exploitation phase. Rule 100200
fires on every \texttt{HONEYPOT\_HIT} event, which is logged whenever
\texttt{/api/admin/config} or \texttt{/api/admin/metrics} is accessed.
Rule 100201 fires when 3 or more honeypot hits arrive from the same IP within
5 minutes, confirming deliberate probing rather than accidental access.

\begin{lstlisting}[style=xmlstyle, caption={Wazuh rules for APT honeypot detection.}]
<rule id="100200" level="10">
  <if_sid>86600</if_sid>
  <field name="event_type">HONEYPOT_HIT</field>
  <description>Honeypot endpoint hit on zone
    aggregator - possible APT probe</description>
  <group>apt,honeypot,</group>
</rule>

<rule id="100201" level="14" frequency="3" timeframe="300">
  <if_matched_sid>100200</if_matched_sid>
  <same_field>client_ip</same_field>
  <description>APT confirmed: 3+ honeypot hits
    in 5 minutes from same IP</description>
  <group>apt,attack,persistent,</group>
  <options>alert_by_email</options>
</rule>
\end{lstlisting}

Rule 100203 intentionally does not trigger an active response. Recon activity
alone (probing unknown paths) is insufficient evidence for a block. The same
pattern could arise from misconfigured clients. Only rule 100201, which confirms
sustained deliberate probing of endpoints with no legitimate callers, activates
the response chain.

\subsection{Response}

Rule 100201 triggers the active response chain described in
Section~\ref{sec:monitoring}. The source IP is blocked at the
\texttt{DOCKER-USER} chain within approximately 5 seconds of the third honeypot
hit. Because the block is applied before any lateral movement to the alert
service can complete, the stolen \texttt{INTERNAL\_TOKEN} cannot be used for
further exploitation from the same IP.

\subsection{Evaluation}

Both phases of the attack are detected correctly. In testing, rule 100203 fires
at approximately 10 seconds into the demo (after 5 probes of unknown paths) and
rule 100201 fires at approximately 35 seconds (after the 3rd honeypot hit). The
IP is blocked before the alert-service pivot begins, as confirmed by the absence
of any rule 100204 (lateral movement) alert (the absence of that alert itself
demonstrates that the defence succeeded). The Grafana \emph{Attack Chain} timeseries
shows a recon burst followed by the honeypot staircase, with the block applied
between the two phases.

The principal limitation is coverage: the honeypot captures only the two
pre-configured administrative paths. A novel misconfigured endpoint not included
in the honeypot list would not generate \texttt{HONEYPOT\_HIT} events and would
therefore not be detected by rule 100201. Additionally, a low-rate recon sweep
of fewer than 5 unknown paths in 30 seconds would not trigger rule 100203, allowing
silent enumeration to proceed.
